\section{范式讨论}
\label{sec:paradigm}

\subsection{智能计算核心：从硅到基底}

前文建立了ICA类比框架及其分层体系结构。图~\ref{fig:paradigm_timeline}总结了经典CPU演进与LLM系统发展之间的结构性对应关系。本节审视核心本身：作为智能系统中\textit{处理器}的大语言模型，以及半个世纪的CPU演进如何照亮其发展轨迹。论证沿六个阶段展开：核心是什么，它如何扩展，异构设计如何绕过扩展极限，核心周围有哪些专用单元，什么基底可能最终取代硅，以及如何衡量它是否正常工作。

\begin{figure}[t]
\centering
\resizebox{\textwidth}{!}{%
\begin{tikzpicture}[
    font=\sffamily,
    >=Stealth,
    % ---- Node styles ----
    cpuNode/.style={
        draw=classical!80, fill=classical!18,
        rounded corners=3pt,
        minimum width=22mm, minimum height=9mm,
        inner sep=2pt, align=center,
        font=\scriptsize\bfseries, text=classical!90!black,
        line width=0.6pt
    },
    llmNode/.style={
        draw=modelnative!80, fill=modelnative!18,
        rounded corners=3pt,
        minimum width=22mm, minimum height=9mm,
        inner sep=2pt, align=center,
        font=\scriptsize\bfseries, text=modelnative!90!black,
        line width=0.6pt
    },
    futureNode/.style={
        draw=modelnative!60, fill=modelnative!8,
        rounded corners=3pt,
        minimum width=22mm, minimum height=9mm,
        inner sep=2pt, align=center,
        font=\scriptsize\bfseries\itshape, text=modelnative!70!black,
        line width=0.6pt, dashed
    },
    yearLabel/.style={
        font=\tiny\bfseries, text=timelineGray!85!black
    },
    timelineLine/.style={
        line width=1.6pt, draw=#1!50
    },
    phaseBand/.style={
        fill=#1, fill opacity=0.08, rounded corners=4pt
    },
    curvedArrow/.style={
        -{Stealth[length=2.2mm, width=1.5mm]},
        thick, draw=#1,
        shorten <=2pt, shorten >=2pt
    },
    arrowLabel/.style={
        font=\tiny\bfseries, text=#1,
        fill=white, fill opacity=0.85, text opacity=1,
        rounded corners=1.5pt, inner xsep=3pt, inner ysep=1.2pt,
        draw=#1!30, line width=0.3pt
    },
    timelineHeader/.style={
        font=\small\bfseries, text=#1,
        fill=#1!6, rounded corners=3pt, inner sep=4pt
    }
]

% ==============================================================
%  COORDINATE SYSTEM
%  x-axis: logical year position (evenly spaced for clarity)
%    CPU:  x = 0, 3, 6, 9, 12, 15, 18
%    LLM:  x = 0, 3, 6, 6, 9, 9, 12   (some share year positions)
%  y-axis:
%    CPU timeline at y = 2.8
%    LLM timeline at y = -1.2
% ==============================================================
\def\cpuY{2.8}
\def\llmY{-1.2}

% ==============================================================
%  PHASE BACKGROUND BANDS
% ==============================================================
% CPU Phase 1: Sequential / Scalar era (1945-2004)
\fill[phaseBand=classical]
    (-1.3, \cpuY+1.1) rectangle (4.7, \cpuY-0.5);
% CPU Phase 2: Parallelism & Heterogeneity (2005-2019)
\fill[phaseBand=classical]
    (4.9, \cpuY+1.1) rectangle (13.7, \cpuY-0.5);
% CPU Phase 3: Specialization (2020+)
\fill[phaseBand=classical]
    (13.9, \cpuY+1.1) rectangle (19.5, \cpuY-0.5);

% LLM Phase 1: Pretrained Foundation (2020)
\fill[phaseBand=modelnative]
    (-1.3, \llmY+0.65) rectangle (1.7, \llmY-1.05);
% LLM Phase 2: Interfaces & Bottlenecks (2023-2024)
\fill[phaseBand=modelnative]
    (1.9, \llmY+0.65) rectangle (7.7, \llmY-1.05);
% LLM Phase 3: Heterogeneous Orchestration (2025-Future)
\fill[phaseBand=modelnative]
    (7.9, \llmY+0.65) rectangle (13.5, \llmY-1.05);

% Phase labels (subtle, inside bands)
\node[font=\tiny\itshape, text=classical!45]
    at (1.7, \cpuY+0.85) {Sequential Era};
\node[font=\tiny\itshape, text=classical!45]
    at (9.3, \cpuY+0.85) {Parallelism \& Heterogeneity};
\node[font=\tiny\itshape, text=classical!45]
    at (16.7, \cpuY+0.85) {Specialization};

\node[font=\tiny\itshape, text=modelnative!45]
    at (0.2, \llmY-0.7) {Foundation};
\node[font=\tiny\itshape, text=modelnative!45]
    at (4.8, \llmY-0.7) {Interfaces \& Bottlenecks};
\node[font=\tiny\itshape, text=modelnative!45]
    at (10.7, \llmY-0.7) {Orchestration \& Specialization};

% ==============================================================
%  TIMELINE HEADERS
% ==============================================================
\node[timelineHeader=classical] at (9, \cpuY+1.55) {CPU Architecture};
\node[timelineHeader=modelnative] at (6, \llmY+1.05) {LLM Systems};

% ==============================================================
%  TIMELINE LINES
% ==============================================================
\draw[timelineLine=classical] (-1.0, \cpuY) -- (19.0, \cpuY);
\draw[timelineLine=modelnative] (-1.0, \llmY) -- (13.0, \llmY);

% ==============================================================
%  CPU TIMELINE MILESTONES  (top row)
% ==============================================================
% 1945: ENIAC
\node[cpuNode] (cpu1) at (0, \cpuY) {ENIAC};
\node[yearLabel, below=2.5mm of cpu1] {1945};

% 1978: x86 ISA
\node[cpuNode] (cpu2) at (3, \cpuY) {x86 ISA};
\node[yearLabel, below=2.5mm of cpu2] {1978};

% 2005: Dennard Scaling Collapse
\node[cpuNode] (cpu3) at (6, \cpuY) {Dennard\\[-1pt]Collapse};
\node[yearLabel, below=2.5mm of cpu3] {2005};

% 2006: Multi-core
\node[cpuNode] (cpu4) at (9, \cpuY) {Multi-core};
\node[yearLabel, below=2.5mm of cpu4] {2006};

% 2012: big.LITTLE
\node[cpuNode] (cpu5) at (12, \cpuY) {big.LITTLE};
\node[yearLabel, below=2.5mm of cpu5] {2012};

% 2016: Heterogeneous SoC
\node[cpuNode] (cpu6) at (15, \cpuY) {Hetero.\\[-1pt]SoC};
\node[yearLabel, below=2.5mm of cpu6] {2016};

% 2020+: Domain-Specific Accelerators
\node[cpuNode] (cpu7) at (18, \cpuY) {Domain\\[-1pt]Accelerators};
\node[yearLabel, below=2.5mm of cpu7] {2020+};

% ==============================================================
%  LLM TIMELINE MILESTONES  (bottom row)
% ==============================================================
% 2020: Pretrained LLMs
\node[llmNode] (llm1) at (0, \llmY) {GPT-3};
\node[yearLabel, above=2.5mm of llm1] {2020};

% 2023: Token/Tool Interface
\node[llmNode] (llm2) at (3, \llmY) {ChatGPT\\[-1pt]Plugins};
\node[yearLabel, above=2.5mm of llm2] {2023};

% 2024: Inference Energy Wall
\node[llmNode] (llm3) at (5.5, \llmY) {KV Cache\\[-1pt]Wall};
\node[yearLabel, above=2.5mm of llm3] {2024};

% 2024: Sparse MoE
\node[llmNode] (llm4) at (8, \llmY) {Sparse\\[-1pt]MoE};
\node[yearLabel, above=2.5mm of llm4] {2024};

% 2025: Heterogeneous Model Orchestration
\node[llmNode] (llm5) at (10.5, \llmY) {Multi-Model\\[-1pt]Orchestr.};
\node[yearLabel, above=2.5mm of llm5] {2025};

% 2025: Specialized Small Models
\node[llmNode] (llm6) at (13, \llmY) {Specialized\\[-1pt]SLMs};
\node[yearLabel, above=2.5mm of llm6] {2025};

% Future: Alternative Substrates  (not connected by arrows; shown for completeness)
% (Omitted to keep the figure focused on the five structural parallels)

% ==============================================================
%  CURVED CONNECTING ARROWS  (CPU → LLM structural parallels)
% ==============================================================
% The arrows curve downward from the CPU timeline to the LLM timeline.
% We use bend left/right to create aesthetically curved paths.

% 1. Dennard Collapse → KV Cache Wall  ("Power Wall")
\draw[curvedArrow=red!70!black]
    (cpu3.south) to[bend right=18]
    node[arrowLabel=red!70!black, pos=0.45, anchor=west, xshift=1pt]
        {Power Wall}
    (llm3.north);

% 2. Multi-core → Sparse MoE  ("Parallelism Shift")
\draw[curvedArrow=blue!70!black]
    (cpu4.south) to[bend right=22]
    node[arrowLabel=blue!70!black, pos=0.45, anchor=west, xshift=1pt]
        {Parallelism Shift}
    (llm4.north);

% 3. big.LITTLE → Multi-Model Orchestration  ("Heterogeneity")
\draw[curvedArrow=orange!85!black]
    (cpu5.south) to[bend right=22]
    node[arrowLabel=orange!85!black, pos=0.45, anchor=west, xshift=1pt]
        {Heterogeneity}
    (llm5.north);

% 4. x86 ISA → ChatGPT Plugins  ("ISA Abstraction")
\draw[curvedArrow=classical!85!black]
    (cpu2.south) to[bend right=12]
    node[arrowLabel=classical!85!black, pos=0.50, anchor=east, xshift=-1pt]
        {ISA Abstraction}
    (llm2.north);

% 5. Domain Accelerators → Specialized SLMs  ("Specialization")
\draw[curvedArrow=purple!70!black]
    (cpu7.south) to[bend right=28]
    node[arrowLabel=purple!70!black, pos=0.40, anchor=west, xshift=1pt]
        {Specialization}
    (llm6.north);

\end{tikzpicture}}% end resizebox
\caption{Structural parallels between CPU and LLM system evolution.
Curved arrows connect milestones that share the same architectural driver:
Dennard scaling collapse and the inference energy wall both arise from a \emph{Power Wall};
the multi-core and MoE transitions reflect a \emph{Parallelism Shift};
big.LITTLE and heterogeneous model orchestration embody \emph{Heterogeneity};
x86 ISA and the token/tool interface represent \emph{ISA Abstraction};
and domain accelerators and specialized small models pursue \emph{Specialization}.}
\label{fig:paradigm_timeline}
\end{figure}


\subsubsection{LLM作为处理器：计算核心为何不需要记忆}

指令集架构（Instruction Set Architecture, ISA）抽象是计算机体系结构中最具深远影响的理念之一。正如Patterson和Hennessy在数十年的学术著作中所论证的，CPU的职责是\textit{忠实地执行指令}，而非存储知识\cite{patterson_hennessy_2017}。同一个程序无论运行在Intel Core i9还是AMD Ryzen上，都编译为相同的ISA；底层硬件的演进独立于上层软件。这种解耦，在RISC运动中得以形式化，并在 successive x86世代中不断精炼，催生了整个现代软件生态。

当前的大语言模型尚未内化这一教训。主流预训练范式将模型视为百科全书：摄入互联网，将知识压缩为权重，然后凭记忆回答问题。这实质上是ENIAC范式：为特定应用场景硬接线的专用硬件\cite{patterson_hennessy_2017}。正如CPU从专用逻辑（ENIAC，1945年）演化为实现稳定ISA的通用处理器（x86，1978年），LLM也必须从``无所不知''的预训练转向``做一个聪明的处理器''，在编排调度下忠实地变换输入。

token接口、上下文窗口和工具使用协议共同构成了初生的\textbf{智能ISA}：模型计算核心（ICA L1，物理执行层）与OS/Agent层（L2，推理服务层）之间的稳定契约。在这一契约下，核心的职责是可靠的指令执行：给定结构化的提示词，生成忠实的响应。知识检索、长期记忆和多步规划属于更高层次，正如文件I/O和进程调度属于OS而非CPU。

这种解耦具有直接的实际收益。安全机制、反越狱防御和常识护栏可以迁移到OS层：上下文过滤器、安全分类器和输出验证器包裹在核心外围，而非内嵌于其权重之中。这一类比是精确的：随着体系结构的成熟，内存保护和特权环从硬件层迁移到了操作系统层\cite{ritchie_thompson_1974}。一个精简的模型核心，被编排层安全机制所包围，既比一个试图内化所有约束的单体模型更可审计，也更加强大。

ISA类比还赋予了体系结构的多样性。正如x86运行在截然不同的Intel和AMD微架构上，同一个Agent框架也可以运行在Transformer、RWKV、Mamba或混合核心上，只要每个核心都实现了智能计算契约\cite{hennessy_patterson_2019}。竞争的焦点从``谁的模型知道更多''转变为``谁的处理器更快、更便宜、更忠实''。ICA框架在L1层容纳这种多样性，同时维持统一的L2层接口。

\subsubsection{频率扩展与思考模式：功耗墙的重演}

1974年，Robert Dennard及其同事发表了一条统治CPU设计长达三十年的扩展定律：随着晶体管缩小，功率密度保持恒定，从而可以在不超出热限制的前提下实现持续提升的时钟频率\cite{dennard_1974}。当Dennard扩展定律在2005年前后失效时，业界撞上了\textbf{功耗墙（power wall）}。频率停滞在约4 GHz附近，唯一的出路是转向多核设计，在每个时钟周期完成更多工作，而非追求每秒更多时钟周期\cite{esmaeilzadeh_2011}。

LLM推理面临类似的\textbf{能耗墙（energy wall）}。token生成成本随参数量线性增长；服务一个405B参数的模型需要数据中心级的电力。``把它做大就行''的时代正在走向终结，原因与GHz竞赛终结的原因相同：物理定律和经济学对单个计算元件能耗散多少能量施加了硬性限制。

两个领域的产业回应遵循着平行路径。Intel的Turbo Boost在面对突发的高负载时自动提升时钟频率，随后在常规负载时降频回撤。LLM的\textbf{思考模式（thinking mode）}（扩展的思维链、扩展推理）运作方式完全相同：系统自动地在困难任务上消耗更多token，从而消耗更多计算资源，同时在简单任务上节约资源。以最大token消耗进行的扩展推理是\textit{超频（overclocking）}的类比：以不成比例的计算成本换取边际质量提升，用效率换取峰值吞吐量。正如持续超频会损毁硬件，持续的最大思考模式在经济上是不可持续的。

\textbf{暗硅（dark silicon）}问题，由Esmaeilzadeh等人（2011年）形式化，直接映射到LLM部署场景。在一块拥有数十亿晶体管的芯片上，能耗预算只允许其中一小部分同时处于活跃状态，其余的则是``暗的''。在一个单体LLM中，在服务延迟和成本约束下，并非所有参数都能为每个token``点亮''。这催生了稀疏激活和混合专家（Mixture-of-Experts, MoE）架构，其中每个输入只激活相关的参数子集，这是多核转型的智力后裔\cite{moore_1965, amdahl_1967}。

历史教训是清晰的。正如功耗墙迫使CPU产业从GHz竞赛转向异构多核设计，推理能耗墙也将迫使LLM系统从单体模型扩展转向异构模型编排。下一节将考察这一转型如何已经在进行中。

\subsubsection{big.LITTLE智能：异构模型编排}

2011年，ARM推出了big.LITTLE架构：在同一芯片上搭配高性能Cortex-A15和能效优异的Cortex-A7，两者都实现相同的ARM ISA\cite{biglittle_arm_2012}。其核心洞见在于，大多数工作负载并不需要峰值性能；一个小巧高效的核心处理日常任务，而强大的核心仅在需求旺盛时激活。调度策略经历了三代演进：\textit{集群迁移（cluster migration）}（只有big\textit{或}LITTLE处于活跃状态）、\textit{HMP}（所有核心同时活跃，由OS级任务分配调度）和\textit{EAS}（能耗感知调度，将硬件遥测集成到Linux内核调度器中，内核5.0，2019年）\cite{eas_linux_2019}。

这三阶段演进直接映射到Agent框架的成熟路径。单模型LLM系统对应集群迁移：一个模型处理一切。规划与执行路由，即由大模型生成计划、小模型执行每一步，对应HMP：多个模型同时活跃，进行粗粒度任务分配。成本感知的动态模型选择，即编排器将每个子任务路由到在可接受质量下最小化预期成本的模型，对应EAS：能耗感知的、遥测驱动的调度\cite{kumar_2003_singleISA, saez_2022_thread_director}。图~\ref{fig:biglittle}展示了这一三阶段对应关系。

% ---------------------------------------------------------------------------
% big.LITTLE  ↔  Heterogeneous LLM Orchestration — three-stage comparison
% fig:biglittle
% ---------------------------------------------------------------------------
\begin{figure}[htbp]
\centering
\providecolor{bigColor}{HTML}{C44E52}
\providecolor{littleColor}{HTML}{4A78A8}
\providecolor{schedColor}{HTML}{3A8A3A}
\providecolor{stageHdr}{HTML}{555555}
%
\resizebox{0.98\textwidth}{!}{%
\begin{tikzpicture}[
    font=\small, >=Latex,
    stagebox/.style={draw=stageHdr!35, fill=stageHdr!4, rounded corners=7pt,
        minimum width=5.8cm, line width=0.6pt},
    bigbox/.style={draw=bigColor!65!black, fill=bigColor!14, rounded corners=4pt,
        minimum width=2.0cm, minimum height=0.80cm,
        align=center, font=\scriptsize\bfseries, text=bigColor!80!black, line width=0.6pt},
    litbox/.style={draw=littleColor!65!black, fill=littleColor!14, rounded corners=4pt,
        minimum width=1.8cm, minimum height=0.70cm,
        align=center, font=\scriptsize\bfseries, text=littleColor!80!black, line width=0.6pt},
    schbox/.style={draw=schedColor!65!black, fill=schedColor!10, rounded corners=4pt,
        minimum width=2.8cm, minimum height=0.80cm,
        align=center, font=\scriptsize, text=schedColor!75!black, line width=0.55pt},
    divline/.style={draw=black!18, dashed, line width=0.55pt},
    darr/.style={-{Latex[length=1.6mm]}, black!45, line width=0.55pt},
    earr/.style={-{Latex[length=2.2mm]}, black!35, line width=0.75pt},
    ann/.style={font=\scriptsize\itshape, text=black!50, align=center},
    hdr/.style={font=\small\bfseries, text=stageHdr!75!black},
    subhdr/.style={font=\scriptsize, text=black!45, align=center},
]

% ── column x-centres ──────────────────────────────────────────
\def\xA{-7.5}   % Stage 1
\def\xB{0.0}    % Stage 2
\def\xC{7.5}    % Stage 3

% ── row y-centres (top section) ───────────────────────────────
% Classical rows
\def\yTitle{4.6}
\def\ySub{4.0}
\def\yCoreTop{2.9}   % top big/LITTLE row
\def\yCoreBot{1.8}   % bottom big/LITTLE row (Stage 2/3)
\def\ySched{1.4}     % scheduler row
\def\yAnnotC{1.9}    % classical annotation (moved up)
% divider
\def\yDiv{0.3}
\def\yDivLbl{0.15}
% LLM rows
\def\yLLMTop{-0.6}
\def\yLLMSched{-2.0}
\def\yAnnotL{-2.6}

% ── stage backgrounds ─────────────────────────────────────────
\node[stagebox, minimum height=8.6cm] at (\xA, 0.55) {};
\node[stagebox, minimum height=8.6cm] at (\xB, 0.55) {};
\node[stagebox, minimum height=8.6cm] at (\xC, 0.55) {};

% ── headers ───────────────────────────────────────────────────
\node[hdr]    at (\xA, \yTitle) {Stage 1: Cluster Migration};
\node[subhdr] at (\xA, \ySub)   {(ARM big.LITTLE, 2011)};
\node[hdr]    at (\xB, \yTitle) {Stage 2: HMP};
\node[subhdr] at (\xB, \ySub)   {(All cores active, OS-level routing)};
\node[hdr]    at (\xC, \yTitle) {Stage 3: EAS};
\node[subhdr] at (\xC, \ySub)   {(Linux 5.0, 2019; telemetry-driven)};

% ══════════════════════════════════════════════════════════════
% STAGE 1 — Classical half
% ══════════════════════════════════════════════════════════════
\node[bigbox] (A_big)    at (\xA-1.7, \yCoreTop-0.5) {big\\[-1pt]Cortex-A15};
\node[litbox] (A_little) at (\xA+1.7, \yCoreTop-0.5) {LITTLE\\[-1pt]A7};
% X mark between the two boxes
\draw[red!55, line width=1.3pt]
    (\xA-0.45, \yCoreTop-0.90) -- (\xA+0.45, \yCoreTop-0.10);
\draw[red!55, line width=1.3pt]
    (\xA-0.45, \yCoreTop-0.10) -- (\xA+0.45, \yCoreTop-0.90);
\node[ann] at (\xA, \yAnnotC-0.5) {\textit{Only one cluster active}\\[-2pt]\textit{at a time}};

% ══════════════════════════════════════════════════════════════
% STAGE 2 — Classical half
% ══════════════════════════════════════════════════════════════
\node[bigbox] (B_big)   at (\xB-1.5, \yCoreTop) {big core};
\node[litbox] (B_lit)   at (\xB+1.5, \yCoreTop) {LITTLE};
\node[schbox] (B_sched) at (\xB,     \ySched)    {OS scheduler};
\draw[darr] (B_sched.north) -- ++(0,0.2) -| (B_big.south);
\draw[darr] (B_sched.north) -- ++(0,0.2) -| (B_lit.south);

% ══════════════════════════════════════════════════════════════
% STAGE 3 — Classical half
% ══════════════════════════════════════════════════════════════
\node[bigbox] (C_big)   at (\xC-1.5, \yCoreTop) {big core};
\node[litbox] (C_lit)   at (\xC+1.5, \yCoreTop) {LITTLE};
\node[schbox, minimum width=3.0cm] (C_sched) at (\xC, \ySched) {Energy Model\\[-2pt]+ telemetry};
\draw[darr] (C_sched.north) -- ++(0,0.2) -| (C_big.south);
\draw[darr] (C_sched.north) -- ++(0,0.2) -| (C_lit.south);

% ── dividers ──────────────────────────────────────────────────
\foreach \cx in {\xA, \xB, \xC}{
    \draw[divline] (\cx-2.6, \yDiv) -- (\cx+2.6, \yDiv);
    \node[font=\tiny, text=black!32] at (\cx, \yDivLbl)
        {Classical\,$\uparrow$\quad LLM\,$\downarrow$};
}

% ══════════════════════════════════════════════════════════════
% STAGE 1 — LLM half
% ══════════════════════════════════════════════════════════════
\node[bigbox] at (\xA-1.3, \yLLMTop-0.5) {Large LLM\\[-1pt]{\tiny(e.g., GPT-4)}};
\node[litbox] at (\xA+1.3, \yLLMTop-0.5) {SLM\\[-1pt]{\tiny(e.g., Haiku)}};
\node[ann]    at (\xA, \yAnnotL+0.3) {\textit{Single model handles}\\[-2pt]\textit{all tasks}};

% ══════════════════════════════════════════════════════════════
% STAGE 2 — LLM half
% ══════════════════════════════════════════════════════════════
\node[bigbox] (B_llm)   at (\xB-1.7, \yLLMTop)   {Large LLM};
\node[litbox] (B_slm)   at (\xB+1.7, \yLLMTop)   {SLM};
\node[schbox] (B_router)at (\xB,     \yLLMSched)  {Plan/execute router};
\draw[darr] (B_router.north) -- ++(0,0.2) -| (B_llm.south);
\draw[darr] (B_router.north) -- ++(0,0.2) -| (B_slm.south);
\node[ann] at (\xB, \yAnnotL-0.3) {\textit{Large model plans,}\\[-2pt]\textit{small model executes}};

% ══════════════════════════════════════════════════════════════
% STAGE 3 — LLM half
% ══════════════════════════════════════════════════════════════
\node[bigbox] (C_llm)   at (\xC-1.7, \yLLMTop)   {Large LLM};
\node[litbox] (C_slm)   at (\xC+1.7, \yLLMTop)   {SLM};
\node[schbox, minimum width=3.0cm] (C_router) at (\xC, \yLLMSched) {Cost-Capability\\[-2pt]Model};
\draw[darr] (C_router.north) -- ++(0,0.2) -| (C_llm.south);
\draw[darr] (C_router.north) -- ++(0,0.2) -| (C_slm.south);
\node[ann] at (\xC, \yAnnotL-0.3) {\textit{Per-subtask cost-optimal}\\[-2pt]\textit{routing (Claude Code, Codex)}};

% ── evolution arrows between stages ───────────────────────────
\draw[earr] (-4.5, \yCoreTop)
    -- node[above, font=\scriptsize, text=black!40] {+concurrency}
    (-2.9, \yCoreTop);
\draw[earr] (3, \yCoreTop)
    -- node[above, font=\scriptsize, text=black!40] {+telemetry}
    (4.5, \yCoreTop);

\end{tikzpicture}}%
\caption{Three-stage evolution of heterogeneous scheduling, mapped between classical ARM big.LITTLE (top half of each panel) and LLM model orchestration (bottom half). Stage~1 (cluster migration / single model) activates only one resource type at a time. Stage~2 (HMP / plan-execute routing) runs both simultaneously with coarse-grained task assignment. Stage~3 (EAS / cost-capability routing) uses real-time telemetry to route each subtask to the cost-optimal compute substrate, the pattern already deployed in Claude Code and Codex.}
\label{fig:biglittle}
\end{figure}


生产级Agent系统已经在实践这一模式。Claude Code在同一Agent循环中路由Sonnet级模型进行规划、Haiku级模型进行执行\cite{liu_2025_claude_code}。Devin使用多Agent编排配合动态规划检查点，随着任务复杂度的变化在模型层级间切换\cite{cognition_2025_multiagent}。这些不是原型系统；它们是服务于数百万查询的已部署系统。

推测解码（speculative decoding）将这一类比延伸到推理服务层。在EAGLE和Medusa等方案中，一个小型草稿模型提出token，由大模型验证，在无质量退化的前提下实现2--3$\times$加速\cite{li_2024_eagle}。这是异构调度的软件类比：小模型（LITTLE核心）处理常规token生成，大模型（big核心）仅在验证或纠正时介入。这一类比甚至延伸到失败模式：正如EAS在过载时回退到SMP调度，Agent系统在小模型能力不足时也会从成本优化的小模型回退到昂贵的大模型。

混合专家（MoE）架构代表了异构计算的\textit{模型内}类比。DeepSeekMoE将专家细分为细粒度单元，每个token稀疏激活部分子集\cite{deepseek_2024_moe}。HMoE明确使用不同大小的专家：较小的专家处理简单模式，较大的专家应对复杂问题，在参数层面实现了异构计算\cite{wang_2024_hmoe}。这就是单一模型内部的big.LITTLE：门控网络充当Thread Director，将token路由到大小合适的专家。

\subsubsection{专用协处理器：小型专用模型的崛起}

现代片上系统（System-on-Chip, SoC）将大量硅面积分配给专用协处理器：GPU用于并行图形处理，NPU用于神经网络推理，ISP用于图像信号处理，DSP用于音频处理。通用CPU编排这些单元，但自身很少执行具体的领域专用工作。同样的模式正在智能系统中浮现，小型专用模型处理那些不需要（也不应承担）完整LLM调用开销的任务。

\textbf{安全模型}是智能栈的安全协处理器。专门为有害内容检测训练的专用内容审核分类器，以极低的成本达到了GPT-4级别的安全性能\cite{harmbench_2024, airbench_2024}。这些模型是TPM模块和硬件安全飞地的类比：用于威胁检测的专用参数，与通用核心隔离，可独立审计和更新。在ICA框架中，它们运行在L2（推理服务层），作为L1计算核心周围的护栏。

\textbf{上下文处理模型}用于检索增强生成，对应经典系统中的缓存层次。交叉编码器排序模型和双通道检索模型管理哪些上下文到达LLM核心，正如带有专用管理硬件的L1/L2/L3缓存减少主存访问\cite{slm_edge_2025}。检索模型是缓存控制器；向量数据库是主存；LLM核心是CPU。一个调校良好的检索流水线可以将大模型的上下文窗口需求降低一个数量级，这正是缓存层次相对于平铺内存访问所提供的效率增益。

\textbf{小型语言模型（Small Language Model, SLM）}（参数量在3B以下）是边缘计算的类比：可部署在设备端，具有可接受的延迟，处理常规推理而无需云端往返\cite{llama32_2024, phi3_2024}。在big.LITTLE框架中，它们是LITTLE核心：Cortex-A7的等价物，处理后台分类、格式转换、实体提取和其他确定性任务。它们的定义特征不是``几乎和大型模型一样好''，而是它们是特定任务的\textit{合适工具}，将大型模型释放给真正需要其能力的任务。

体系结构蓝图是Apple M1 SoC：4P+4E CPU核心、集成GPU、16核神经网络引擎，全部共享统一内存\cite{apple_m1_2020}。LLM系统的等价物是一个异构集成：大型规划模型、小型执行模型、安全分类器、检索排序器和专用分类器，全部在统一编排下共享公共上下文池。ICA L1层承载这种多样性；L2提供使集成作为连贯系统运作的调度和内存管理。

\subsubsection{超越硅：类脑、生物与量子基底}

Hennessy和Patterson的图灵奖演讲（2019年）呼吁以领域专用架构作为突破Dennard扩展定律终结和摩尔定律极限的路径\cite{hennessy_patterson_2019}。同样的迫切性驱动着对非冯诺依曼基底用于智能计算的探索。这不是推测性的未来主义；这是ISA论证所预测和要求的基底演进。

\textbf{类脑硬件（Neuromorphic Hardware）}已达到部署规模。Intel的Hala Point系统集成了1,152颗Loihi 2芯片，实现了11.5亿个神经元，采用事件驱动的脉冲计算，在适合的工作负载上实现了100--1000$\times$的能效提升\cite{hala_point_2024}。SpiNNaker2以脑启发拓扑部署了约500万个ARM核心\cite{spinnaker2_2024}。早期工作已在Loihi 2芯片上演示了LLM推理，包括无矩阵乘法的语言模型\cite{abreu_2025_loihi2_llm}。这些并非GPU的通用替代品，而是面向脉冲和事件驱动工作负载的领域专用加速器，这类工作负载正是生物神经计算所偏爱的。

\textbf{生物计算（Biological Computing）}已从研究走向产品。Cortical Labs的CL1于2025年3月出货，是首台商用生物计算机：活体神经元在微电极阵列上培养，与数字系统交互\cite{cortical_2025_cl1}。普林斯顿大学的3D-MIND将70,000个活体神经元与可编程柔性电子器件集成\cite{3dmind_2025}。FlyWire联盟完成的果蝇全脑连接组（Nature，2024年）使Eon Systems在整脑仿真中实现了95\%的行为准确率\cite{flywire_2024}。类器官智能（Organoid Intelligence）路线图将当前概念验证系统通往实用生物处理器的路径正式化\cite{hartung_2023_oi}。

\textbf{存算一体（Compute-in-Memory）}架构通过消除冯诺依曼瓶颈实现了显著的效率提升。面向LLM注意力机制的模拟存内计算展示了相比数字实现100$\times$的延迟降低和10,000$\times$的能耗降低\cite{leroux_2025_cim}。Taalas HC1芯片走向了逻辑极端：将Llama 3.1 8B的权重永久硬连线到掩模ROM中，实现了17,000 tokens/sec的吞吐量，完全无需动态内存访问\cite{taalas_2026_hc1}。

\textbf{量子计算（Quantum Computing）}用于LLM仍处于早期阶段。Quixer（Quantinuum，2024年）是首个在真实硬件上运行的量子原生Transformer\cite{khatri_2024_quixer}。理论工作探索了面向Transformer推理的量子线性代数\cite{guo_2024_quantum_transformer}，但容错的、应用级的量子计算预计还需要5--15年以上。量子和经典的混合方法是近期的可行路径。

这些多样的基底强化了第1小节的ISA论证。一个稳定的智能计算接口，即``智能ISA''，允许物理基底从硅GPU演进到类脑芯片，再到生物神经组织，再到量子处理器，而不会破坏软件栈。正如x86 ISA使软件免受从真空管到晶体管到集成电路的转变影响，token/工具使用协议使Agent框架免受从GPU上的Transformer权重到Loihi芯片上的脉冲模式再到CL1阵列上活体神经元的转变影响。ICA L1层抽象了这种基底多样性；L2及以上层永远不需要知道是什么物理介质在执行它们的指令。

\subsubsection{基准测试核心：木桶原理与最坏情况评估}

经典CPU基准测试经历了一段痛苦的成熟过程。在1990年代，``兆赫兹营销''让消费者相信时钟频率是处理器质量的唯一衡量标准。实际情况更为复杂：流水线深度、缓存大小、分支预测准确率和指令级并行度都影响着真实性能。业界最终收敛到SPEC和Geekbench等综合套件，测量多样工作负载而非峰值时钟频率\cite{spec_benchmark}。

LLM基准测试面临同样的成熟挑战。MMLU平均分和排行榜得分是智能系统的``兆赫兹''：单一数字掩盖的远比揭示的多。牛津大学对445个LLM基准测试的审查发现，84\%缺乏基本的统计检验\cite{fodor_2025_benchmarks}。在标准测试中得分90\%的模型，在未见问题上可能仅得2\%，这正是MHz神话的精确类比，时钟频率是真实性能的误导性代理指标\cite{benchmark_pitfalls_2025}。

\textbf{最弱环节定律（Law of the Weakest Link）}，在ICLR 2025上形式化，为\textbf{木桶原理（bucket principle）}提供了理论基础：一个系统的智能受限于其最弱的能力维度，而非平均值\cite{zhong_2024_weakest_link}。一个在事实回忆上卓越但在逻辑推理上失败的模型，或者一个在良性查询上表现优美但在对抗探测下产生危险输出的模型，不是一个``总体还不错''的系统，而是一个存在关键失效模式的系统。木桶从最低点漏水。

最坏情况指标正在涌现以应对这一问题。\textit{Worst-at-k}量化多次采样响应中的预期最坏情况性能，捕获尾部风险而非平均行为\cite{worst_at_k_2025}。安全基准测试如HarmBench和AIR-BENCH评估系统最危险的失效模式而非平均情况的帮助性\cite{harmbench_2024, airbench_2024}。这些是LLM领域对应于实时系统中最坏情况执行时间（WCET）分析和硬件验证中对抗性压力测试的等价物。

对于智能计算核心本身，基准测试必须评估处理器对其ISA契约的\textit{忠实度}。LLM核心是否可靠地执行Agent的指令？是否遵循指定的输出格式？是否尊重工具使用协议？这些问题衡量的是处理器的可靠性，而非百科全书式的知识。一个在MMLU上得分70\%但忠实执行99.9\%结构化指令的核心，优于一个在MMLU上得分90\%但频繁偏离其编程规范的核心。ICA框架明确区分了这一点：知识评估属于记忆和检索层；执行可靠性属于计算核心。

CPU基准测试的历史轨迹，从兆赫兹到SPEC再到工作负载特定的性能剖析，表明LLM评估也将沿着同样路径发展：从聚合分数到逐能力维度的性能剖析，再到与真实Agent任务挂钩的工作负载专用基准。木桶原理加速了这一转变，因为它证明了从业者已经直觉到的：平均情况指标对于单点失效可能造成灾难性后果的系统来说是不够的。





\subsection{智能操作系统：从代码到世界}

\subsubsection{历史弧线：从任务专用系统到通用操作系统}

早期的操作系统并非通用系统。它们是为特定使命而构建的工具：IBM 704上的FORTRAN Monitor System（FORTRAN监控系统）服务于科学计算；Colossus时代的系统服务于军事加密；早期的批处理系统服务于薪资处理。没有人会把那时的系统描述为``通用''的操作环境。它们是窄领域、面向特定任务、恰好管理硬件资源的工具。这一特征精确映射到今天的智能Agent系统上。Codex、Claude Code和Devin虽然功能强大，但它们的应用范围局限于软件工程工作流。它们编辑代码仓库、运行终端、管理文件系统、调用API。从历史视角看，它们正是智能计算领域的FORTRAN Monitor System：任务专用工具，日后将被视为更通用系统的先驱。

经典操作系统并非一夜之间变得通用。它需要经历一系列能力拐点：多道程序设计（multiprogramming，Atlas，1962年）使并发任务执行成为可能；分时系统（time-sharing，CTSS，1961年；Multics，1965年）使交互式计算成为可能；虚拟内存与保护模式（Unix、VMS，1970年代）提供了隔离与可靠性；网络协议栈（BSD TCP/IP，1980年代）连接了不同系统；硬件抽象层（Windows NT HAL，1990年代）将操作系统与特定硬件解耦 \cite{ritchie_thompson_1974, corbato_1962_ctss, corbato_1965_multics}。每一层能力的叠加都是必要的，此后一个操作系统才能同时服务于科学计算、办公自动化、实时控制和多媒体。

当前的智能系统处于其演进过程中的``编程语言阶段''。正如早期操作系统为特定硬件运行汇编或单语言程序一样，今天的Agent操作系统为特定领域运行``提示程序''，主要是代码。要跃迁到通用智能操作系统，即能够跨机器人、自动驾驶、工程、农业、娱乐和太空导航编排Agent的系统，需要与当年将Unix从单语言、单硬件前驱中分离出来的同类抽象层突破。

上节记录的五代框架追溯了从批处理（第I代）到受治理操作系统（第V代）的演进弧线。但这一轨迹指向第V代之后，即所谓的第VI代：与完整物理和数字世界交互的通用智能操作系统。Patterson和Hennessy指出，ISA是指令集架构（Instruction Set Architecture）的缩写，是实现软硬件协同演进的奠基性抽象 \cite{patterson_hennessy_2017}。Hennessy和Patterson在2019年图灵奖演讲中将开放ISA和领域特定协同设计视为``新黄金时代''的支柱 \cite{hennessy_patterson_2019}。同样的结构性要求，即层间开放且明确定义的接口，将主导从今天仅面向代码的Agent操作系统到明天覆盖全球的智能操作系统的转变。

\subsubsection{计算基底无关性：ISA解耦原则的推广}

上一小节已梳理了类脑、生物、存算一体和量子等非硅基底的最新进展。本节回答一个更根本的架构问题：\textbf{为什么智能操作系统应该与计算核心的物理实现无关？}

答案直接来自ISA抽象的教训 \cite{patterson_hennessy_2017}。ISA之所以是经典计算中最重要的抽象，不是因为某一种实现更好，而是因为它允许软件生态和硬件生态\textbf{独立演进}。ICA的接口契约（第~\ref{sec:icam-interfaces}节）将这一原则推广到模型原生计算：如果一个计算核心---无论是硅基GPU、神经形态芯片、存算一体基底、还是生物神经网络---能够达到定义的智能评分阈值（下一小节讨论），它就应该能够运行智能操作系统。操作系统与计算核心完全解耦，正如Linux与CPU通过ISA解耦一样。

其哲学延伸遵循RISC-V的开放ISA理念 \cite{hennessy_patterson_2019}：任何满足规格的硬件都可以运行软件，无论其物理实现方式如何。甚至一个经过训练的生物网络---一个通过智能评分阈值的鼠脑类器官---原则上也可以作为智能操作系统的计算核心。今天在NVIDIA H100上调度Claude Code子Agent的同一编排层，明天可以调度Cortical Labs CL1上的感觉运动任务 \cite{cortical_2025_cl1}，只要两个核心满足相同的接口契约。这就是最完整形式的\textbf{基底无关性}（substrate agnosticism）：操作系统不关心其计算核心是蚀刻在硅片上还是在培养皿中培养出来的。

\subsubsection{智能评分作为ISA合规机制}

上小节已建立了木桶原理和最坏情况评估的理论基础。本节将视角从``如何评估核心''上移到``如何用评分执行ISA契约''。

ISA合规测试是执行软硬件契约的机制。一颗芯片不能声称实现了x86，除非它通过了覆盖每条指令、每个特权模式、每个边界条件的全面验证套件。没有这种测试，ISA契约就毫无意义。直接类比之下，智能评分标准是智能ISA的执行机制。没有它，基底无关性（上一小节）将沦为一堆无法验证的声明。

Zhong等人在ICLR 2025上发表的``最弱环节定律''为合规的严格性提供了理论基础 \cite{zhong_2024_weakest_link}：当多种能力必须组合才能产生可靠行为时，整体系统性能受限于最弱的单项能力。一个在代码生成上得分95百分位但在安全合规上仅得20百分位的系统，其系统智能为20百分位。木桶中最低的那块板决定了水位。

这意味着\textbf{没有基于木桶原理的评分标准，基底无关性是危险的}。一个在10个安全维度中通过9个但在第10个上失败的生物计算核心，在基于平均值的指标下可能被部署为``合规''系统，造成灾难性风险。最低分执行不仅仅是一种测试偏好，而是一项安全上的刚性要求。正如Blem等人（2013）证明微架构比ISA本身更重要 \cite{blem_2013_microarchitecture}，智能评分必须评估实现质量，而不仅仅是接口合规。

\subsubsection{OS级异构调度：从能耗模型到能力-代价模型}

上小节已建立了异构模型编排的微架构类比。本节将视角上移一层：当异构核心由智能操作系统统一调度时，\textbf{调度策略}本身如何设计？

经典操作系统为这一问题提供了完整的架构蓝图。能耗感知调度（Energy Aware Scheduling, EAS）于2019年合入Linux内核5.0，为操作系统调度器提供了能耗模型（Energy Model），可以预测将每个任务放置到每个CPU上的代价 \cite{eas_linux_2019}。智能操作系统的类比是一个\textbf{能力-代价模型}（Capability-Cost Model），预测将每个子任务路由到每个可用计算核心的质量、延迟和美元成本，无论是大模型、小模型、神经形态芯片还是生物基底，从而实现智能最优调度。

两个关键的OS级调度模式值得关注：

\textbf{过载回退}：当系统过载时，EAS回退到传统的SMP负载均衡，牺牲能效以换取吞吐。智能操作系统的类比是``不惜代价保质量''（quality-at-all-costs）的回退策略：当Agent系统面临峰值负载时，无论成本优化如何，都将任务路由到更大、更昂贵的模型。

\textbf{实时遥测与动态路由}：Intel的Thread Director（Alder Lake，2021年）向操作系统调度器提供关于线程行为的实时硬件遥测数据 \cite{intel_thread_director_2021}。智能操作系统的类比是一个\textbf{任务导向器}（Task Director），实时监控子任务特征（推理深度、安全敏感度、延迟容忍度），并在执行过程中动态重新路由到最优的计算基底。

引入生物和神经形态基底后，调度挑战在质上变得更加困难。生物核心可能具有与硅基GPU根本不同的延迟特征、故障模式和能力签名。智能操作系统调度器不仅要应对``快与慢''的区别，还要应对``概率式与确定式''``高带宽与低延迟''以及``硅基与生物''的区别。这是经典操作系统设计者从未面对过的复杂性层次的异构调度，但其架构模式---能耗模型、过载回退和实时遥测---直接继承了big.LITTLE的设计手册。

\subsubsection{物理世界接口：从纯代码到具身智能}

今天的智能操作系统只与一个领域交互：软件工程。Codex、Claude Code和Devin操作代码仓库、终端、文件系统和API。在上一小节建立的历史框架中，这是智能的FORTRAN Monitor System阶段：面向单一世界的任务专用操作系统。下一个拐点是能够与物理世界交互的智能操作系统：机器人、自动驾驶、工程仿真、农业、娱乐和太空导航。

全脑仿真里程碑正在使具身智能（embodied intelligence）变得具体可行。FlyWire联盟绘制了果蝇约140,000个神经元和数百万个突触的连接图谱 \cite{flywire_2024}，使Eon Systems能够创建行为准确度达95\%的具身虚拟果蝇 \cite{eon_systems_2025_virtual_fly}。在小鼠全皮层仿真（900万神经元、260亿突触）在Fugaku超级计算机上的运行展示了向更复杂具身智能体演进的轨迹 \cite{allen_riken_2025_mouse_fugaku}。这些不是概念验证，而是模拟完整生物体与虚拟环境交互的可用系统。

计算基底无关的智能操作系统意味着，今天调度Claude Code子Agent的同一编排层，原则上明天可以调度机器人操作任务、农业监控Agent和自动驾驶规划器。其要求有两方面：计算核心必须达到智能评分阈值（上一小节讨论的木桶原理），操作系统必须具备适当的世界接口驱动程序：传感器、执行器、仿真环境和遥测数据流。

这直接类比了经典操作系统如何从批处理扩展到实时控制。Unix最初是一个文本处理系统，但随着实时调度、设备驱动程序和网络协议栈的加入，它得以运行工厂自动化、电信交换机和航天器制导系统。每一个新世界需要新的驱动接口，但不需要新的操作系统内核。同一个Linux内核可以运行Web服务器、工厂机器人和火星车。

ICA的L4（语义接口层）和L5（编排层）正是为这种可扩展性而设计的。工具模式和Agent协议就是智能操作系统的设备驱动程序。添加机械臂驱动或卫星遥测接口遵循与今天添加MCP工具服务器相同的抽象模式。接口契约是一致的：感知、规划、执行、验证，无论目标是代码仓库还是玉米地。

生物计算基底为具身智能增添了独特的维度。生物神经网络可能在感觉运动任务、模式识别、自适应控制和实时响应方面具有硅基系统所不具备的天然优势。Kagan等人（2022）证明了生物神经元可以通过反馈学习玩Pong \cite{kagan_2022_pong_bioneurons}，证实了具身生物计算的可行性。智能操作系统应该能够将感觉运动子任务分派到生物核心，同时将符号推理路由到硅基GPU，这是应用于具身认知的异构调度。调度器为正确的任务选择正确的核心，正如经典操作系统将浮点运算分派到FPU、将整数运算分派到ALU。物理世界不是另一个操作系统，而是同一操作系统的另一组驱动程序。


%% ============================================================
